\documentclass[11pt,a4paper]{article}
\usepackage[margin=26mm]{geometry}
\usepackage{amsmath,amssymb,amsthm,mathtools}
\usepackage{lmodern,microtype,booktabs,longtable,array}
\usepackage[colorlinks=true,linkcolor=blue,citecolor=blue,urlcolor=blue]{hyperref}
\hypersetup{pdftitle={A cubic zero-free boundary from plain-moment parameter feedback},pdfauthor={RH-Weil research project}}
\newtheorem{theorem}{Theorem}[section]
\newtheorem{proposition}[theorem]{Proposition}
\newtheorem{lemma}[theorem]{Lemma}
\newtheorem{corollary}[theorem]{Corollary}
\theoremstyle{definition}
\newtheorem{definition}[theorem]{Definition}
\newtheorem{remark}[theorem]{Remark}
\newcommand{\Rea}{\operatorname{Re}}
\newcommand{\eps}{\varepsilon}
\newcommand{\bet}{\beta_*}
\newcommand{\sig}{\sigma_*}
\newcommand{\pos}[1]{\left(#1\right)_+}
\newcommand{\src}[1]{\nolinkurl{#1}}
\makeatletter
\g@addto@macro\UrlBreaks{\do\a\do\b\do\c\do\d\do\e\do\f\do\g\do\h\do\i\do\j\do\k\do\l\do\m\do\n\do\o\do\p\do\q\do\r\do\s\do\t\do\u\do\v\do\w\do\x\do\y\do\z\do\0\do\1\do\2\do\3\do\4\do\5\do\6\do\7\do\8\do\9}
\makeatother
\title{A cubic zero-free boundary\\from plain-moment parameter feedback}
\author{RH--Weil research project}
\date{October 7, 2026\\[3pt]\small A derivation relative to explicitly imported results}
\begin{document}
\maketitle
\begin{abstract}
We extend the usable parameter range of the generic plain moment
in the September 30, 2026 preprint \emph{The Quasi-Riemann Hypothesis},
by replaying its proof from its specified underlying inputs.
The range $37/50\le\kappa\le1$ retains the original moving conductors,
natural zero extensions, coefficient classes and uniform mesh.
Coupling the actual moment parameter to the whole-family zero
supremum yields the strict zero-free half-plane
\[
 \Rea s>\sigma_*=0.874957019420098946\ldots,
 \qquad
 7884\sigma_*^3-18819\sigma_*^2+14643\sigma_*-3686=0,
\]
with the root specified by an exact rational isolating interval.
The conclusion is for finite-order Hecke functions over
$\mathbb Q(\sqrt{-3})$ and, by the imported quadratic transfer,
all Dirichlet functions, including zeta. We prove a continuous
quadratic certificate at a reference parameter and pay the cost of
using the larger actual parameter. We retain the complete low Gram
estimate and verify strict detector widths, every physical row range,
Euler domains, the principal residue and uniform family continuation.
This improves the preceding free-$b$ boundary by approximately
$5.04\cdot10^{-8}$. The theorem assumes an explicit imported package
and its Hecke seven-eighths bootstrap; it is not an independent
certification of those inputs or of an external proof kernel.
The principal pole is allowed and the boundary line is excluded.
No new critical-line zero proportion is asserted.
\end{abstract}
\begingroup
\small
\tableofcontents
\endgroup

\section{Statement, source and hypotheses}
\label{sec:inputs}
Let $F=\mathbb Q(\sqrt{-3})$. We use the ideal and element
representatives, norm $q_a$, units, ray conditions and zero-extension
conventions of the pinned source \cite{OAI}. A superscript $S$ deletes
the finite Euler factors in the final fixed arithmetic datum.
This paper is separate from the papers giving $69999/80000$ and
the free-$b$ boundary \cite{Previous,FreeB}; those delivered artifacts
are preserved. The new step is a moment-domain extension and
parameter feedback, rather than another choice of $b$ in the old
$\kappa=3/4$ envelope.

\begin{definition}[Imported package $\mathcal R$]\label{def:R}
We assume the following source results, with their complete coefficient
classes, masks, uniformities and side conditions:
finite Gauss and full-correlation identities, reciprocity, complete
support and natural zero extensions; coefficientwise probe and Poisson
identities, finite compensation, local arithmetic tables and ray
calibration; smooth separation, Gaussian annuli, reflection sectors,
generic reflected energy and the full quantitative additive Gram bound;
global logarithmic and buffered reciprocal control, primitive functional
equations, growth and complete contour tails; detector witnesses,
marked and recursive inverse moments and sixth-power amplification;
fixed-field counting and fixed-ray prime asymptotics; the whole
finite-order Hecke seven-eighths zero-free theorem; and the quadratic
factorization transferring that theorem to Dirichlet functions.
The new plain-moment range below is derived from the underlying
identities and estimates, not assumed by extending the statement of
the original plain theorem outside its published range.
\end{definition}

Define
\[
 \bet=\sup\bigl(\{1/2\}\cup
 \{\Rea\rho:L_F(\rho,\eta)=0,\ \rho\text{ nontrivial},
                  \ \eta\text{ primitive of finite order}\}\bigr).
\]
The principal pole is excluded. The only bootstrap used in the
proof of the new boundary is $\bet\le7/8$ from $\mathcal R$.
The previous free-$b$ boundary is used only for comparison.
Pure norm twists change imaginary parts, not this real supremum.

The source is frozen at commit
\src{adc7f1241b42e322a6451854ab7e4b4c146bf78a}.
Its September 30 \src{build/paper.tex} has canonical LF SHA-256
\src{42a5ee0febca59fd1def55cfd6c6808c322ef4237d7726303ea52f711deac6a3}.
Line ranges below refer to that file. Table~\ref{tab:R} specifies
the interfaces used, so that the relative nature of the result
is explicit.

\begin{longtable}{@{}>{\raggedright\arraybackslash}p{.27\textwidth}>{\raggedright\arraybackslash}p{.23\textwidth}>{\raggedright\arraybackslash}p{.43\textwidth}@{}}
\caption{Imported interfaces and retained conditions.}\label{tab:R}\\
\toprule Interface & Source lines & Use in this paper\\ \midrule
\endfirsthead
\toprule Interface & Source lines & Use in this paper\\ \midrule
\endhead
Smooth, Gaussian and global log control & 1123--1209, 1288--1315,
1536--1594 & Joint density, all heights; the new real strip is compact
inside the original global domain.\\
Actual detector and marked moments & 4220--4685, 9220--9269 &
Original presentation, witnesses, coefficient group and strict widths.\\
Generic plain proof & 12492--14974 &
Underlying proof replayed in Section~\ref{sec:plain}; no out-of-range
invocation of its original conclusion.\\
Physical slots and counts & 14985--15446, 15920--15992 &
Whole-slot mean gain, conjugation, masks, inducing exceptions and
shared mesh; new crossing derived below.\\
Reflection and Gram & 7663--7804, 8340--8364 &
All frozen parts, actual dyads and complete third Gram term.\\
Physical identities and local bounds & 8707--9049, 9101--9195 &
Same finite compensation, normally convergent corrections and full
completed sums.\\
Whole-bin, principal and outer & 5808--5896, 6021--6463 &
Global move before local partition; new parameter domains verified
in Section~\ref{sec:contours}.\\
Uniform choice and continuation & 16191--16463, 6520--6580,
400--502 & Losses and mesh before slots; target before internal
height orders; external tails last.\\
Dirichlet transfer & 6705--6780 &
Quadratic factorization, finite Euler deletions and principal pole.\\
\bottomrule
\end{longtable}

Let $e$ be the unique root in $(1/6,167/1000)$ of
\begin{equation}\label{eq:cubic}
 p(e)=657e^3-954e^2+21e+20=0.
\end{equation}
It is strictly isolated by
\begin{equation}\label{eq:bracket}
 \frac{16683858898627}{10^{14}}<e<
 \frac{16683858898628}{10^{14}}.
\end{equation}
Set
\begin{equation}\label{eq:parameters}
 \sig=\frac{11}{12}-\frac e4,\qquad
 \kappa_*=\frac56-\frac e2=2\sig-1,\qquad
 b=\frac{-5181+156335e-387630e^2}{81941}.
\end{equation}

\begin{theorem}\label{thm:main}
Assume $\mathcal R$. Every finite-order Hecke $L$-function over $F$
has no zero in $\Rea s>\sig$, with the root specified by
\eqref{eq:cubic}--\eqref{eq:parameters}.
The same holds for all Dirichlet $L$-functions, including $\zeta$.
Principal functions may have their pole at $1$.
The boundary line is not included.
\end{theorem}

Substitution gives the cubic for $\sig$ in the abstract.
The numerical values
\[
\begin{aligned}
 e&=0.166838588986270882\ldots,\\
 \kappa_*&=0.749914038840197892\ldots,\\
 b&=0.123405685419541002\ldots .
\end{aligned}
\]
are for orientation; the proof uses rational arithmetic and
the specified real embedding. We make no claim of optimality
over all possible probes or arithmetic estimates.

\section{Replaying the plain moment below three quarters}
\label{sec:plain}
The original generic plain statement imposes $\kappa\ge3/4$.
It is essential to reprove the statement in the enlarged range.
We describe every numerical dependency on that lower endpoint;
the source identities remain imported inputs.

\begin{proposition}[Extended plain range]\label{prop:plain}
Use the source normalized plain sums $S_{\psi_k}(n)$, positive
prime-slot product $Q_{\psi_k}$ and row family $\mathcal R_z$.
Retain rows $q_k\ll Z^m$, the complete displayed moving radical
$Z^q$ of the fixed twist, and effective width $M=m+q$.
Retain the same polynomial-size fixed extra puncture mask,
the finite coefficient group $\Theta$, its fixed finite-ray
coefficient lists, disjoint underlying prime supports and
nonprincipal inducing-family exclusions.
For each $\eps>0$ and fixed bounded length range, choose the
numerical mesh $\eta_{\rm mesh}>0$ as specified below.
Then, for $37/50\le\kappa\le1$, positive-slot data satisfying
\[
\begin{gathered}
 \max_i z_i\le\eta_{\rm mesh},\qquad n_1+n_2+6\kappa z\le M,\\
 \bet\le(1+\kappa)/2\quad(\kappa<1)
\end{gathered}
\]
obey
\begin{equation}\label{eq:plain}
 \sum_{k\in\mathcal R_z}
 |S_{\psi_k}(n_1)S_{\psi_k}(n_2)Q_{\psi_k}|^2
 \ll Z^{M+\eps},
\end{equation}
with the original finite seminorm and polynomial-height factors.
For $z=0$ retain the original bounded nonnegative plain lengths.
The numerical mesh can be chosen after the bounded length range and
$\eps$, before the slot count, uniformly in this $\kappa$ interval
and all admissible moving moduli and outer labels.
\end{proposition}
\begin{proof}
Write $A=n_1+n_2+z$, so the positive-slot budget is
$A+(6\kappa-1)z\le M$. Natural zeros remain present even when a
phase cancels or is six-divisible; changing a phase does not erase
a deleted prime from the moving radical.

The prime estimate in source lines 12833--12900 is obtained on
the displaced line to the right of
$s_\kappa=(1+\kappa)/2\ge\bet$.
Now $s_\kappa\ge87/100$, within a fixed compact real strip of
the global logarithmic estimate. For the same coefficient expansion
and nonprincipal exclusions this gives
\[
 |Q_{\boldsymbol\omega}|^2
 \le C Z^{\kappa z+\eps_1}p_j(\boldsymbol W)^a
                         (1+|\boldsymbol\omega|)^t .
\]
The contour displacement is the original $2\eps z$; height
and derivative orders are paid by the same weighted Fourier measure.
At $\kappa=1$ the absolute prime count suffices.
There is no new principal residue.

At the terminal-width step (12925--12938),
\[
 z\le M/(6\kappa)\le25M/111,\qquad \kappa z\le M/6.
\]
The original displayed $z\le2M/9$ is replaced, but it has no later
use in the proof. The terminal exponent uses only $\kappa z\le M/6$.

For the same-band comparison $A>5M/6$ (12947--12980),
\[
 (6\kappa-1)z<M/6,\qquad
 6\kappa-1\ge86/25,\qquad z<25M/516<M/20.
\]
The two comparison bounds consequently remain
\[
 A_{\rm comp}\le3M/2-A+2z+\xi\le23M/30+\xi,
\]
\[
 A_{\rm comp}+(6\kappa-1)z
 \le5M/2-2A+2z+\xi\le14M/15+\xi .
\]
Thus the original $M/15$ margin and $\xi\le\rho/30$ still suffice.
The strict inequality $25/516<1/20$ has slack $1/645$.

The aggregate support boxes, whole-dyad tails, first complete
extraction, two M\"obius indicators, common characters, $p^6$ pool,
second correlation and child-width decrease in 13095--13875
do not use the old lower endpoint. Their actual child width is
$M'=M+J-g-g_2+t_2\le M-\sigma_{\rm drop}$.
The original pool separation
$\ell_{\rm pool}=\sigma_{\rm drop}/3$ and
$\eta_{\rm pool}<\sigma_{\rm drop}/6$ is retained.

In the actual clipping and mask-erasure steps (13876--14310),
shortening or removing slots does not increase the affine
budget because $6\kappa-1\ge0$. The actual defect remains
\[
 F_{\rm act}
 =(A_{\rm act}-M'_{\rm act}+(6\kappa-1)z_{\rm act})_+
 \le6(w+\ell_{\rm edge})+\delta_{{\rm fr},1}+2\theta_N.
\]
The whole-slot removal cost is
$\kappa d_z\le F_{\rm act}/6+\kappa\eta_{\rm pool}
             \le F_{\rm act}/6+\eta_{\rm pool}$.
Separate actual rectangles, paired shells, zero extensions and
all moving conductor factors are retained; there is no extra
cost from extending the interval.

The inducing-$\Theta$ exceptions, equal-product cancellation,
common centering, sixth-power counts, padding and reflection in
14312--14777 are independent of this lower endpoint.
The positive-slot condition still gives $A\le M$.
Finally the numerical Lipschitz bound in 14779--14974 requires
only $0\le6\kappa-1\le5$. Keep the original width floor, bands,
finite depth and strict drop $\sigma_{\rm drop}/2$.
Within each band perform uncentered before centered estimates;
establish all zero-slot cases before the positive-slot induction.
The same finite induction closes with the same terminal cost and
comparison margins.

All numerical losses, mesh and support thresholds are chosen
before the fixed slot count. The numerical constants above are
uniform in the compact $\kappa$ interval. After fixing slot and
profile data, choose the finite internal seminorm and height orders
in the original reverse induction order; an external tail order
chosen later cannot alter them. This proves the original uniform
contract \eqref{eq:plain} from $\mathcal R$ in the stated range.
\end{proof}

The source's full support, correlation and analytic estimates
are still assumptions. The finite algebra audit in the appendix
does not verify this induction or those imported estimates.

\section{Actual inverse and plain counts}
\label{sec:counts}
Put $\alpha=5/6$. In nonfloor bins use
$\delta=2a-1\in[1/50,3/4]$ and the source length-weighted
mean amplitude $q$ over \emph{all} physical slots.
Error and zero slots have gain zero. Let
$x=q/\delta\in[0,1/2]$ and $y=1/2-x$.
For $37/50\le\kappa\le1$ define
\begin{equation}\label{eq:DP}
 c=(3\kappa)^{-1},\quad B_{\rm pl}=2-2cx,\quad
 D=3-(1+2c)x,\quad P=B_{\rm pl}(1-x).
\end{equation}
The actual witness has
$r+m\ge t-o(1)$, $r\le t+o(1)$ and $m\le1/2+o(1)$,
with its original inverse and plain spikes.
Use two copies of the same plain witness; its fourth spike is
not replaced by a single squared spike. The resulting count lines are
\[
 A_I(r)=1-\delta\{x+(1-x)r\},\qquad
 S_t(r)=1-\delta\{cx+B_{\rm pl}(t-r)\}.
\]
Their crossing is
\[
 r_*(t)=\frac{B_{\rm pl}t-(1-c)x}{D},\qquad
 R_{\rm short}(t)=1-\delta+\frac{\delta P}{D}(3/2-t).
\]
The long branch remains
$1-\delta+(\alpha-\delta)(t-1)$.
Equating the short and long branches gives
\begin{equation}\label{eq:R}
 J=(\alpha-\delta)D+\delta P,\qquad
 t_\kappa=1+\frac{\delta P}{2J},\qquad
 R_{*,\kappa}=1-\delta+
           \frac{(\alpha-\delta)\delta P}{2J}.
\end{equation}
These formulas preserve the same character and presentation.
Witness and physical heights need not coincide; their parameter
Sobolev factors and cumulative $T_1/2$ allowance are retained.

For later strict-width checks, the rectangle satisfies
\[
 D\ge455/222>2,\quad P\ge86/111>3/4,\quad
 35/48\le J\le5/2,\quad 1<t_\kappa<3/2,
\]
\[
 r_*(1)\ge8/13>3/5,\quad r_*(3/2)=1,\quad
 1/3\le t-r_*(t)\le1/2\quad(1\le t\le3/2).
\]
The inverse request has $z_M=(1-r)/2\le5/26<1/5$;
the plain request has $z_P=(1-2m)/(6\kappa)\le25/333<1/5$.
On the inverse branch request $z\le z_M-\nu_0$, for a fixed
$\nu_0>0$, and retain both marked widths:
\[
 1-r-2z\ge2\nu_0,\qquad
 3-2r-8z
 =4(1-r-2z)+(2r-1)\ge8\nu_0+1/5-o(1).
\]
On the plain branch request $z\le z_P-\nu_0$, giving
$1-2m-6\kappa z\ge6\kappa\nu_0$.
Capacity-zero neighborhoods, $r\ge1$ and $m\ge1/2$ use the
original no-slot or amplified endpoints.
Thus a physical supply $e/d>1/5$ with fixed margin suffices.
Whole-slot rounding, $\Theta$ expansion and conjugation of the
witness and all selected factors together remain essential.

\section{The same physical probe and the complete low estimate}
\label{sec:low}
Set
\begin{equation}\label{eq:geometry}
 l_x=(1-e-b)/2,\quad l_y=(1-e+b)/2,\quad
 M=1-e,\quad h=(1+3e+b)/2,\qquad X=Z^{l_x},\ Y=Z^{l_y}.
\end{equation}
Choose fixed even $K$, equal $\ell_i=e/K$, $P_i=Z^{\ell_i}$
and the original disjoint nonnegative smooth prime windows.
Let $I_\eta$ be the source completed physical probe, and
$I_{\eta;D}$ that same expression with $1_{D\mid cn^3}$ inserted.
For $p_J=\prod_{i\in J}p_i$, $p_\varnothing=1$, retain precisely
\[
\begin{split}
 I_{\eta,\rm mod}(Z)
 =\sum_{(p_i)}\prod_i W_i(q_{p_i}/P_i)
 \sum_{J\subseteq\{1,\ldots,K\}}&
 (-1)^{|J|}q_{p_J}^{-3/2}\overline{\eta(p_{J^c})}\\
 &\cdot I_{\eta;p_{J^c}}
 (X/q_{p_J},Y/q_{p_J},Zq_{p_{J^c}}).
\end{split}
\]
All ray data, annular norm ratios, masks and windows are unchanged.
The completed $cn^3$ sums are infinite and absolutely convergent
under $\mathcal R$; only the compensation subsets are finite.
The principal Mellin exponent is
\begin{equation}\label{eq:C}
 C_b(s)=s+l_x/2-1+h/6=s-2/3-b/6 .
\end{equation}
Here the raw kernel's $z$ coefficient is $h-e$, and the exact
tuple adds $e/6$ at the residue. The combined coefficient is
$h/6$; it is not charged a second $e/6$.

\begin{proposition}\label{prop:low}
For the parameters \eqref{eq:parameters}, under the generic low
inputs of $\mathcal R$,
\[
 |I_{\eta,\rm mod}(Z)|
 \ll_{\eta,\rm fixed,\eps}Z^{(1-e)/4-b/6+\eps}
 =Z^{C_b(\sig)+\eps}.
\]
\end{proposition}
\begin{proof}
We replay the parameter-independent reflected-energy and Gram
calculation of \cite{FreeB}. For a compensation subset write
$d_J=\sum_{i\in J}\ell_i$, $e'=e-d_J$, $M'=M-2d_J$.
The exact tuple ratio lies in a fixed positive compact interval;
thus $X'\asymp Z^{l_x-d_J}$, $Y'\asymp Z^{l_y-d_J}$
and the common row ball is $Q\asymp Z^{M'}$.
Surviving marked slots are summed before the row norm;
one common Gaussian profile and Fourier density is retained.

For the actual squarefree dyad $H$ and frozen powerful length $O$,
the sector conditions are
$\Delta_H=M'-O-H\ge-O(1/\log Z)$,
$2A_0\le O+O(1/\log Z)$, $N_0\le A_0$ and $z_a\le e'$.
The dual length is
$T_d=2H+2A_0+2z_a-1-e'-\theta_N-N_0-3B_0$.
In particular $T_d-(H-3d_J)\le O(1/\log Z)+|\theta_N|$.
The generic reflection bound uses
$s_{\rm hyb}=\min\{v,z_a,(v+z_a)/3\}$.
On its $s_{\rm hyb}=z_a$ branch the exponent is at most
$M'-\Delta_H$; on the other branch $s_{\rm hyb}\ge v/2$.
The frozen-part saving is the exact identity
\[
\begin{split}
 O/2+\Delta_H+S_0+B_0+T_d/4
 ={}&\frac{2M'+2z_a-1-e'+2\Delta_H-\theta_N}{4}\\
 &+\frac{2A_0-N_0+B_0+4S_0}{4}.
\end{split}
\]
The second line is nonnegative. Keeping the kernel saving before
squaring and counting frozen parts once gives the complete row estimate
\[
 \sum_{q_m\ll Q}|B_m^J|^2
 \ll Z^{M'+(5e-1+d_J)_+/4+\eps}.
\]
All far Gaussian annuli are first paid by the full absolute input.
No step depends on $\kappa$ or discards a shared prime or empty slot.

The new parameters satisfy $1/6<e<1/5$, $b>0$,
$M'\ge1-3e>0$ and $l_x-d_J,l_y-d_J>0$.
Consequently $Q,Y'\ge1$ eventually in fixed polynomial ranges.
The source parameter $P_a=Y'^2/Q\asymp Z^b$ is eventually at least
$1$. Use the complete Gram bound
\[
 \sum_{q_m\ll Q}|A_m(Y')|^2
 \ll (Q/Y')(1+P_a^{1/6}+P_a^2/Y')Z^\eps .
\]
Cauchy with the original separation factor $Q^{-1/2}$ gives a
fixed tuple exponent
\[
 (X')^{1/2}
 Z^{\frac12\max\{0,b/6,2b-l_y+d_J\}
       +(5e-1+d_J)_+/8+\eps}.
\]
The tuple count, $q_{p_J}^{-3/2}$ and shortened $X'^{1/2}$
together contribute $-d_J$.
Relative to $l_x/2+b/12$, the additional exponent is
\[
 F_{\rm low}(d)=-d+\tfrac12(11b/6-l_y+d)_+
                         +\tfrac18(5e-1+d)_+ .
\]
Its every slope is at most $-3/8$, so the maximum is at $0$.
The rational interval checks below give the stronger all-subset
gap $l_y-e-11b/6>2/25$. Therefore the complete exponent is
$l_x/2+b/12=(1-e)/4-b/6$.
All small reflection, separation, height and logarithmic losses
are absorbed in the prescribed $\eps$ in the original order.
Finally \eqref{eq:parameters} and \eqref{eq:C} give the equality.
\end{proof}

\section{A continuous certificate at the reference parameter}
\label{sec:certificate}
Only in this section set $\kappa=\kappa_*$ as an algebraic reference.
It is not an analytic moment assumption when $\bet>\sig$.
Writing $y=1/2-x$ transforms \eqref{eq:DP} into
\[
 D=(5/2-c)+(1+2c)y,\qquad
 P=(1-c/2)+2y+2cy^2.
\]
Put
\[
 K_0=-1/4+5e/4+b/6,\qquad W=1/2+3e/2-ey,
\]
and define the reference endpoint on the actual rectangle by
\[
 E_*=K_0+(1/2+e)\delta+e\delta x-h(1-R_{*,\kappa_*})
     =K_0+W\delta-h(1-R_{*,\kappa_*}).
\]
Direct expansion, retaining the full whole-slot mean gain, gives
\begin{equation}\label{eq:F}
 -2JE_*=F=A(y)\delta^2-B(y)\delta+C(y),
\end{equation}
where
\begin{align}
 A&=-2(P-D)(W-h)+hP,\notag\\
 B&=2K_0(P-D)+\tfrac53D(W-h)+\tfrac56hP,\label{eq:ABC}\\
 C&=-\tfrac53DK_0.\notag
\end{align}
The cleared polynomial $F$ is defined for arbitrary real $\delta$;
the identity with $E_*$ is used only where $J\ne0$.

Let $Q=4AC-B^2=\sum_{j=0}^4 Q_jy^j$.
For general $e,b,\kappa$, $Q_0$ is a strictly concave quadratic
in $b$, its $b^2$ coefficient being
$-(3312\kappa^2-936\kappa+67)/(1728\kappa^2)$.
The numerator has negative discriminant and positive leading
coefficient. Its maximizer is
\begin{equation}\label{eq:bopt}
 b_{\rm opt}=
 \frac{5040e\kappa^2-900e\kappa+25e-432\kappa^2+36\kappa+3}
 {3312\kappa^2-936\kappa+67}.
\end{equation}
The maximum is
\[
 -\frac{25(6\kappa-1)^2(15\kappa-2)
 \{(1314\kappa-159)e^2-(36\kappa+6)e-30\kappa+5\}}
 {162\kappa^2(3312\kappa^2-936\kappa+67)} .
\]
At $\kappa=5/6-e/2$, the braces are $-p(e)$.
Reducing \eqref{eq:bopt} modulo \eqref{eq:cubic} gives exactly
the $b$ in \eqref{eq:parameters}. Hence $Q_0=0$.

All remaining signs have rational certificates. Write
$A=\sum_{0}^3 A_jy^j$ and $C=C_0+C_1y$.
Table~\ref{tab:signs} gives strict rational enclosures, obtained
by reducing the coefficients in $\mathbb Q[e]/(p)$ and applying
interval Horner to \eqref{eq:bracket}.
The endpoints are rational numbers, not decimal evidence.

\begin{table}[ht]
\centering
\begin{tabular}{@{}cccc@{}}
\toprule Coefficient & Strict enclosure & Coefficient & Strict enclosure\\
\midrule
$A_0$ & $(47/100,48/100)$ & $C_0$ & $(7/100,8/100)$\\
$A_1$ & $(121/100,122/100)$ & $C_1$ & $(6/100,7/100)$\\
$A_2$ & $(86/100,87/100)$ & $Q_1$ & $(4/100,5/100)$\\
$A_3$ & $(29/100,30/100)$ & $Q_2$ & $(19/100,20/100)$\\
 & & $Q_3$ & $(26/100,27/100)$\\
 & & $Q_4$ & $(7/100,8/100)$\\
\bottomrule
\end{tabular}
\caption{Positive coefficient certificate; $Q_0=0$ exactly.}
\label{tab:signs}
\end{table}

\begin{proposition}\label{prop:certificate}
On $1/50\le\delta\le3/4$, $0\le x\le1/2$, the reference
endpoint satisfies $E_*\le0$. Its unique equality is
\[
 y=0,\qquad \delta_*=\frac{5-9e}{6+18e},\qquad
 R_{*,\kappa_*}=2/3 .
\]
\end{proposition}
\begin{proof}
For every $y\ge0$, Table~\ref{tab:signs} implies $A>0$ and
$Q\ge0$, with $Q=0$ only at $y=0$. Thus for arbitrary real
$\delta$ the cleared polynomial obeys
\[
 F=A\left(\delta-\frac B{2A}\right)^2+\frac Q{4A}\ge0 .
\]
On the actual rectangle $J>0$, so \eqref{eq:F} gives $E_*\le0$.
At $y=0$, reduction modulo $p$ gives
$B(0)/(2A(0))=(5-9e)/(6+18e)$ and $R=2/3$.
The same rational bracket gives $1/50<\delta_*<3/4$.
This proves the continuous assertion without inferring it from
finite samples.
\end{proof}

\section{The actual parameter and its feedback cost}
\label{sec:feedback}
Assume for contradiction $\bet>\sig$ and put $\Delta=\bet-\sig>0$.
The bootstrap gives $\Delta\le(e-1/6)/4$.
The analytic calls to Proposition~\ref{prop:plain} must use
\begin{equation}\label{eq:actual}
 \kappa_{\rm act}=2\bet-1\in(\kappa_*,3/4]\subset[37/50,1],
 \qquad \kappa_{\rm act}-\kappa_*=2\Delta .
\end{equation}
Then $\bet=(1+\kappa_{\rm act})/2$ is the true prime-bound
prerequisite. Calling the lemma at $\kappa_*$ instead would
assume the conclusion to be proved.

\begin{lemma}\label{lem:derivative}
Uniformly on the count rectangle and $37/50\le\kappa\le1$,
\[
 0\le\partial_\kappa R_{*,\kappa}
 =\frac{(\alpha-\delta)^2\delta x(1-x)^2}{3\kappa^2J^2}
 \le\frac{625}{36963}<\frac1{50}.
\]
\end{lemma}
\begin{proof}
Differentiate $P/D$ first:
$\partial_c(P/D)=-2x(1-x)^2/D^2$.
Then differentiate \eqref{eq:R} and use
$\partial_\kappa c=-1/(3\kappa^2)$.
For the bound use $J\ge(\alpha-\delta)D$, $D\ge2$,
$\delta\le3/4$ and $x(1-x)^2\le4/27$.
The resulting bound is
$(3/4)(4/27)/(12(37/50)^2)=625/36963$.
\end{proof}

For the same endpoint geometry, comparison to the actual
$C_b(\bet)$ gives exactly
\[
 E_{\bet}(h)=E_*(h)-\Delta+
 h\{R_{*,\kappa_{\rm act}}-R_{*,\kappa_*}\}
 \le-24\Delta/25 .
\]
Here $h<1$. The same-source full selected-row expression is
\begin{equation}\label{eq:fullE}
\begin{split}
 E_{\bet}(d)={}&K_{\bet}+(1/2+e)\delta+eq
             -h(1-R)\\
 &+(d-h)(R+\delta/2-17/50),\quad
 K_{\bet}=2/3+e+b/6-\bet ,
\end{split}
\end{equation}
with $R=R_{*,\kappa_{\rm act}}$.
It is the expansion of
$a-\bet+h(17/50-1/6)-a l_y-e/2+eq
 +d(R+\delta/2-17/50)$.
Strict local labels and numerator estimates consume the same
conductor deficit once; the gain is still $eq$.
On $1/2\le d\le h$ its slope is at least
$33/50-\delta/2\ge57/200$.
Hence the endpoint controls all these physical lengths.
Choose $\zeta=\Delta/32$; the slope is less than $2$,
so $h<d\le h+\zeta$ costs at most $\Delta/16$.
The central saving before other real losses is therefore
\begin{equation}\label{eq:saving}
 (24/25-1/16)\Delta=359\Delta/400.
\end{equation}
The gap $\Delta$ and actual parameter are fixed before the
mesh, slot count or target character.

\section{All physical ranges and strict arithmetic admission}
\label{sec:admission}
The exact rational bracket and parameter definitions give
\begin{align}
 l_y-e-11b/6&>2/25,\quad
 5e-h-\zeta>1/50,\label{eq:geometrymargins}\\
 h+\zeta&<1,\notag\\
 E_{\rm floor}&=-6/25+(32/25)e+b/6<-1/200,\notag\\
 E_{\rm middle}&\le-71/600+(329/400)e-(421/1200)b<-1/50,\notag\\
 E_{\rm small}&=(13/75)h-l_y/2+63/5000<-2/25.\label{eq:ranges}
\end{align}
These inequalities are recalculated at the cubic parameters,
not inherited numerically from \cite{FreeB}.
The bound
$\zeta\le(e-1/6)/128<1/384000$
also follows exactly.

The floor $a=51/100$ uses $\delta=1/50$, $R=1$, $q\le1/100$
without manufacturing a witness. On $1/100\le d\le1/2$ the original
no-slot count is $R=76/75-2\delta/3$.
Its endpoint value is
$-73/300+(13/50)e-(49/300)b
 +\delta(1/6+3e/4-b/4)$.
The coefficient of $\delta$ is positive, so $\delta=3/4$
gives the middle bound in \eqref{eq:ranges}.
The small bound is already relative to $C_b(\bet)$; no second
$\Delta$ is deducted from it.

Equation~\eqref{eq:geometrymargins} gives
$e/(h+\zeta)>1/5$. It exceeds both actual requests in
Section~\ref{sec:counts} by fixed margins.
Only $d\ge1/2$ uses selected primes; each physical length
$w_i=\ell_i/d\le2\ell_i$.
Choose count and moment losses, fixed capacity decrement,
uniform mesh and rounding first; then take even $K$ sufficiently
large that all $w_i$ fit that mesh and the whole-slot rounding
loss is small. Amplitude and bin widths follow.
Retain the original finite $\Theta$ coefficient combinations,
natural row zeros and disjoint physical and internal prime pools.
All factors and witnesses are conjugated together on the plain
side; separated heights retain their common Sobolev costs.

Outside-$S$ ramification gives the original inducing exclusions.
Remaining $S$-supported sixth-power-free rows are finite and
are covered by the bounded-row argument, not by a false generic
moment call. The cumulative $T_1/2$ and all original profile
and height fees are retained.
Thus \eqref{eq:R} is an actual count with strict widths,
not merely a nominal interpolation.

\section{Euler domains, the principal term and outer rows}
\label{sec:contours}
The parameter change does not alter a local coefficient,
zero extension or compensation operation. Retain the raw weight
\[
 X^{1/2-z}Z^{s+z-1}Y^{w-1}e^{(s+z-1)^2}
                 M(z)\widehat W_1(w)
\]
and the complete selected-prime sum with its coefficientwise
$G_p$ factors and normally convergent unselected correction.
There is no division by a general local correction $H_p$.
The quotient is used only in the principal and dynamic domains
where its near-one condition has been proved.

On the principal box
$\Rea s\ge\sig$, $\Rea w\ge19/20$, $\Rea z\ge33/200$,
the original coefficient tables give
\[
\begin{split}
 c_{\rm good}
 &=\min\{6\sig-401/100,\sig-3/50,47/50\}
   =\sig-3/50>81/100,\\
 c_{\rm ram}
 &=\min\{\sig-1/20,3\sig-3/2\}
   =\sig-1/20>82/100 .
\end{split}
\]
All denominators retain their separation from zero.
The general local term $4-6s-6z$ includes
$\theta=(-\Rea w)_+$; in $D_1(\eps_0)$ use
$\eps_H=\min(\eps_0,1/50)$.
These are full normal-convergence and height domains.
Whole-bin contours are moved globally before the original
buffered amplitude and error partitions.

For the principal row retain exactly
\[
 H_\eta(s)=\prod_{p\notin S}H_p(s,1,1/6).
\]
A target-independent $P_0$ from the positive majorant gives
$|H_\eta-1|\le1/2$ for $\Rea s\ge\sig$.
The final target-dependent enlargement of $S$ only decreases
that majorant; it is the same $S$ in all physical and principal
expressions. For the original fixed-ray windows,
\[
 S_i(Z)=\sum_{p\in\mathcal P_i(Z)}
 W_i(q_p/P_i)q_p^{-5/6}\sim c_iP_i^{1/6}/\log P_i,\quad c_i>0,
\]
\[
 A_T=(-1)^K Z^{-e/6}\prod_i S_i(Z),\qquad
 c_S=\widehat W_1(1)M(1/6)
      (\operatorname{Res}_{v=1}\zeta_F^S(v))^2/6>0 .
\]
Thus $A_T$ is eventually nonzero and
$|A_T|\asymp(\log Z)^{-K}$, with subpower inverse.
No new normalizer is introduced.

The global principal move, residues at $w=1,z=1/6$ and
remaining lines $19/20$, $33/200$ give
\[
 m_w=l_y/20>0,\qquad m_z=h/600>0,\qquad
 0<\mu<\sig\min_i\ell_i .
\]
The double-residue $s$ line moves only to the right, to $2$.
The four local principal error exponents remain
$-\Rea s$, $-6\Rea z$, $4-5\Rea s-6\Rea z$ and
$1-\Rea w-6\Rea z$; their weakest decay at the residue is $\sig$.

For small rows $d<1/100$ and bounded nontrivial units, use
the local proof on $D_1(1/3)$ since $\sig+1/2>4/3$;
do not invoke a numerically stronger old $D_1(3/8)$ condition.
The full numerator and absolute tuple give exactly the
small bound \eqref{eq:ranges}.
For $U>Z^{h+\zeta}$ use the full completed absolute sum on
$(2,2,z_\infty)$, whose exponent is
\[
 B_0+(h+\zeta)(1+\eps)-\zeta z_\infty+\eps,\qquad
 B_0=l_x/2+1+l_y=7/4-3e/4+b/4.
\]
After $\zeta$ is fixed, select sufficiently large fixed
$z_\infty$ before the target, then the small losses and
internal height orders. This pays any required outer saving.
An order is not changed with a moving row or with $Z$.

\section{Uniform choices and the family contradiction}
\label{sec:continuation}
Geometry, the bootstrap, $\Delta$ and $\kappa_{\rm act}$ precede
all target choices. Allocate small fixed fractions of $\Delta$
to count and moment losses, capacity decrement, mesh and rounding.
Then fix even $K$ and all windows, followed by
$\mu=\sig e/(2K)>0$, amplitude widths, prime losses and buffers.
Keep total other central costs below $\Delta/4$ in
\eqref{eq:saving}. Floor, middle and small losses retain at
least half their fixed margins; principal losses are separately
smaller than $m_w/4,m_z/4,\mu/4$.
Fix $\zeta,z_\infty$ in this same target-independent phase.
This order avoids choosing the mesh or $K$ backwards from a
$K$-dependent $\mu$.

Set
\[
 m_0=\min\{\Delta,m_w,m_z,\mu,1/200,1/50,2/25\},\qquad m=m_0/4 .
\]
For a target $\eta$, fix its arithmetic datum, the final $S$
and then the finite internal Sobolev, seminorm and height orders.
For all sufficiently large $Z$ define the physical function
and for all positive $Z$ define the principal function by
\[
 J_\eta(Z)=I_{\eta,\rm mod}(Z)/(c_SA_T(Z)),\qquad
 f_\eta(Z)=\frac1{2\pi i}\int_{(2)}
 Z^{C_b(s)}e^{(s-5/6)^2}
             \frac{H_\eta(s)}{L_F^S(s,\eta)}\,ds .
\]
Both are independent of the auxiliary cutoff $T_1$.
Take the total low and inverse-normalizer loss
$\omega=\Delta/2<\Delta$. Then
\[
 |J_\eta(Z)|\ll_\eta Z^{C_b(\sig)+\omega}.
\]
The common principal and all high estimates give finite
$A_\eta,B_\eta$ such that, for an external order $N$,
\begin{equation}\label{eq:contract}
 |J_\eta-f_\eta|\ll_{\eta,N}
 Z^{C_b(\bet)-m}(1+T_1)^{A_\eta}+Z^{B_\eta}T_1^{-N}.
\end{equation}
The internal orders giving $A_\eta,B_\eta$ are fixed before $N$.
For the actual detector ceiling $\tau_{0,\eta}>0$, choose
\[
 0<\tau_\eta\le
 \min\{1/10000,\tau_{0,\eta},m/[4(A_\eta+1)]\},\qquad
 T_1=Z^{\tau_\eta}.
\]
Choose $N_\eta$ afterwards so that
$B_\eta-N_\eta\tau_\eta<C_b(\bet)-m/2$,
and then the eventual threshold in $Z$.
The common high saving is $\sigma_{\rm hi}=m_0/8>0$.
Target-dependent constants, orders and thresholds do not change it.

\begin{proof}[Proof of Theorem~\ref{thm:main}]
The low and high bounds imply
local uniform convergence of
\[
 \mathcal F_\eta(s)=
 \int_0^\infty f_\eta(Z)Z^{-C_b(s)}\,\frac{dZ}{Z}
 \quad\text{on }\Rea s>\bet-\eps_*,
 \qquad \eps_*=\min\{\Delta/2,m_0/8\}>0 .
\]
At small $Z$ use only $f_\eta$, whose Gaussian integral has
rapid decay by rightward shifts. There is no need to define
$J_\eta$ or divide by $A_T$ there. At large $Z$ use
\eqref{eq:contract} and the strict low gap.
Since $C_b$ has slope $1$, the same Mellin/Fourier identification
as in $\mathcal R$, initially on line $2$, gives
\[
 \mathcal F_\eta(s)=e^{(s-5/6)^2}
                    H_\eta(s)/L_F^S(s,\eta).
\]
Analytic continuation makes this reciprocal expression holomorphic
on the larger half-plane, without moving a contour across a
target zero. Since $\bet-\eps_*>\sig$ and $|H_\eta-1|\le1/2$
there, a zero cannot be canceled by the numerator.

The supremum need not be attained. Its definition supplies a
target zero with real part $>\bet-\eps_*$.
All common choices preceded that target, and the holomorphic
reciprocal contradicts its existence. Thus $\bet\le\sig$.
Finite deleted Euler factors have no zero in $\Rea s>0$,
so imprimitive functions have the same strict conclusion.
The quadratic factorization in $\mathcal R$ transfers it
to the Dirichlet family. At $1$ the principal pole is allowed;
the zeta factor uses $L(1,\chi_{-3})=\pi/(3\sqrt3)\ne0$.
\end{proof}

\section{Exact comparison and scope of the advance}
The preceding free-$b$ paper gives
$\sigma_0=(1507-2\sqrt{921})/1653$.
Its associated $e_0$ is the positive root of
$1653e_0^2-66e_0-35=0$, isolated by
\[
 \frac{16683838746}{10^{11}}<e_0<
 \frac{16683838748}{10^{11}}.
\]
The lower endpoint for $e$ in \eqref{eq:bracket} exceeds
this upper endpoint. Hence $e>e_0$, and exactly
$\sig<\sigma_0$.
For orientation,
$\sigma_0-\sig=0.000000050379069793\ldots$.
This comparison is not used as a bootstrap.

The advance is the re-established plain contract below $3/4$
and its actual count feedback. The theorem remains relative
to $\mathcal R$. It does not independently certify the source
paper, its analytic inputs or an external proof kernel, and
does not establish RH. The separate critical-line proportion
route still needs a complete signed response and finite-constant
fourth-moment budget; none is supplied by this strip gain.
The equality of the reference envelope is a statement about a
count certificate, not a constructed bad row or zero.

\appendix
\section{Reproducible exact certificate}
\label{app:audit}
The file \src{scripts/hybrid_kappa_feedback_exact_audit.py}
uses only rational arithmetic. It implements
$\mathbb Q[e]/(p)$ in the basis $(1,e,e^2)$, with
$e^3=(-20-21e+954e^2)/657$, and solves the exact three-dimensional
multiplication system for division.
Modulo $7$, $p(j)$ for $j=0,\ldots,6$ has residues
$6,3,4,3,1,6,5$; the cubic is irreducible.
Endpoint signs and $p'<0$ on $(1/6,167/1000)$ certify the
selected real embedding.

For a direct way to reproduce Table~\ref{tab:signs}, put
$D_0=5/2-c$, $D_1=1+2c$, $P_0=1-c/2$ and
$T_0=P_0-D_0$, $T_1=2-D_1$. Since $W-h=-b/2-ey$,
\begin{align*}
 A_0&=bT_0+hP_0,& A_1&=bT_1+2eT_0+2h,\\
 A_2&=2bc+2eT_1+2hc,& A_3&=4ec,\\
 B_0&=2K_0T_0-5bD_0/6+5hP_0/6,\\
 B_1&=2K_0T_1-5bD_1/6-5eD_0/3+5h/3,\\
 B_2&=4K_0c-5eD_1/3+5hc/3,\\
 C_0&=-5K_0D_0/3,& C_1&=-5K_0D_1/3.
\end{align*}
Then $Q_j=4\sum_{r+s=j}A_rC_s-\sum_{r+s=j}B_rB_s$,
with absent coefficients zero. Reduce these expressions modulo
$p$ and use interval Horner on \eqref{eq:bracket}.
The strict rational enclosures in the table follow; they were
also independently verified with a rational Bernstein enclosure.
The complete rational coefficient triples are saved in
\src{output/hybrid-kappa-feedback-exact-audit.json}.

Run from the repository root:
\begin{center}
 \texttt{python}\ \src{scripts/hybrid_kappa_feedback_exact_audit.py}
\end{center}
The audit has 149 counted checks, including 98 identity checks
at 49 finite endpoint and square-completion models. It also
contains uncounted internal arithmetic assertions. The continuous proof
is coefficient positivity and square completion, not those
49 samples. Its scope is displayed exact algebra, root isolation,
strict parameter inequalities and finite identity models.
It does not verify the plain induction, imported analytic
inputs, actual prime sums, infinite contours, the zero-free
conclusion, RH or an external kernel. Its source-hash field is
fixed metadata: the script does not read or authenticate the external
source or Note 451. Actual version binding is performed separately
by the independent reviews and build manifest.

The lower-$\kappa$ source replay and full research derivation
are preserved in \cite{CountReplay,Note450,Note451}.
Two full independent reviews of this paper, its final source
and PDF hashes, compilation and all-page visual verification
are recorded in
\src{reviews/2026-10-07/kappa-feedback-paper-build-manifest.json}.
The reviewed frozen research draft and the independently
reviewed paper are separate evidence objects.

\clearpage
\begin{thebibliography}{9}
\raggedright
\bibitem{OAI}
OpenAI mathematical preprint,
\emph{The Quasi-Riemann Hypothesis}, September 30, 2026,
commit \src{adc7f1241b42e322a6451854ab7e4b4c146bf78a}.
\url{https://github.com/openai/math/blob/adc7f1241b42e322a6451854ab7e4b4c146bf78a/preprints/The-Quasi-Riemann-Hypothesis-September-30-2026/build/paper.tex}.
\bibitem{Previous}
RH--Weil research project,
\emph{A small improvement of the seven-eighths zero-free boundary
from a compensated probe},
2026. \src{papers/seven-eighths-boundary-improvement-paper.tex}.
\bibitem{FreeB}
RH--Weil research project,
\emph{Optimizing compensated-probe geometry below seven eighths},
2026. \src{papers/free-b-compensated-probe-boundary-paper.tex}.
\bibitem{CountReplay}
RH--Weil research project,
Critical-count witness and lower-$\kappa$ source replay, October 7, 2026.
\src{reviews/2026-10-07/hybrid-critical-count-witness-research.md}.
\bibitem{Note450}
RH--Weil research project, Note 450, October 7, 2026.
\src{notes/450-plain-kappa-extension-and-actual-capacity.md}.
\bibitem{Note451}
RH--Weil research project, Note 451, October 7, 2026.
\src{notes/451-kappa-feedback-cubic-boundary-and-family-continuation.md}.
\end{thebibliography}
\end{document}
